\section{Related work}

Partial-identification analysis characterizes the target values compatible with an observed law \citep{Manski1990,HorowitzManski1995,Tamer2010}. For a known distribution of assignment probabilities and a fixed released trial law, the sharp average treatment effect interval developed below is the bounded-outcome, overlap specialization of \citet[Theorem 3.2 and eq. (3.1)]{FanShermanShum2014}. Within each released label, its endpoints account for the possible association between outcomes and assignment probabilities.

Data-combination and conditional-transport analyses also characterize outcomes consistent with observed marginal or conditional distributions. In particular, \citet[Section 3]{FanParkPassShi2025} obtain convex identified sets and transport-computable support functions for affine moment models with incomplete data. Here the fixed-law interval is one part of a release-design problem: the analysis ranges over compatible released trial laws to measure the largest attainable interval width, then chooses a finite-label rule to minimize that width. The score-based randomization model and this outer optimization are specific to the present problem.

Statistical matching and record-linkage studies analyze information recovered by linking separate files. \Citet[Section 4, especially Theorem 2]{Komarova2018} show that a positive bound on individual disclosure risk generally leaves their model parameter only partially pseudo-identified. Their object is a probabilistic matching rule based on personal identifiers and a disclosure constraint. Our released object is instead a deterministic function of the assignment score attached to every trial row; its information loss is measured by sharp ATE ambiguity, and point identification occurs exactly when the label determines the score almost surely.

Coarsening also appears in treatment-effect estimation with linked covariates. Under unconfoundedness, outcome smoothness, and restrictions on extreme scores, \citet[Assumptions 1--3 and Informal Theorem 1.1]{KalavasisMehrotraZampetakis2024} construct a data-dependent coarse IPW estimator whose error accounts for propensity-score inaccuracy and sampling. \Citet[Sections 3--5]{LeeWeidner2026} obtain treatment-effect bounds and inference by pooling outcomes across small groups of observations with linked covariates, including settings with limited overlap. For identification and known-score design, our experiment retains a deterministic label in place of the row-level score and assumes a known score marginal; the independent-score-log extension learns that marginal from a separate sample. Under a positive density floor, the resulting optimal ambiguity is of order \(K^{-1}\), and honest length in the known-score experiment is of order \(K^{-1}+n^{-1/2}\). These rates quantify label design under a different observation scheme and different regularity conditions.

The inference questions fit the literature on confidence sets for partially identified parameters \citep{ImbensManski2004,Stoye2009,ChernozhukovHongTamer2007,KaidoMolinariStoye2019}. With a known score distribution, the criterion is worst-case expected length for a jointly chosen label rule and uniformly honest interval. With an independent score log, the rule is fixed and the criterion is worst-case expected positive excess length beyond the population sharp interval from estimation with both released trial rows and the independent score log. These two criteria measure, respectively, the length of an honest interval under label coarsening and the additional length generated by the two sampled data sources.
